% Generated draft: requires mathematical and source review.
\documentclass[11pt]{article}
\usepackage[margin=25mm]{geometry}
\usepackage{fontspec}
\setmainfont{Latin Modern Roman}
\usepackage{amsmath,amssymb,mathtools,unicode-math}
\setmathfont{Latin Modern Math}
\usepackage{enumitem,needspace,xurl,hyperref}
\setlength{\parindent}{0pt}
\setlength{\parskip}{5pt}
\setlength{\emergencystretch}{3em}
\newcommand{\N}{\mathbb{N}}
\newcommand{\Z}{\mathbb{Z}}
\newcommand{\Q}{\mathbb{Q}}
\newcommand{\R}{\mathbb{R}}
\newcommand{\C}{\mathbb{C}}
\newcommand{\F}{\mathbb{F}}
\newcommand{\A}{\mathbb{A}}
\newcommand{\PP}{\mathbb{P}}
\newcommand{\E}{\mathbb{E}}
\newcommand{\eps}{\varepsilon}
\newcommand{\dd}{\,\mathrm{d}}
\begin{document}
\section*{UnsolvedMath ID: 11100019}
\textbf{Source problem code:} catalogue entry
\subsection{Definitions and mathematical statement}
Fix a prime $p$ and height $n$. Let $S$ denote the source's Morava stabilizer group and $E_n$ its intended Morava-$E$ theory. In the usual homotopy-fixed-point framework the relevant spectral sequence has the form
\[
 H_c^s(S;(E_n)_t)\Longrightarrow\pi_{t-s}(E_n^{hS}).
\]
Whether $S$ is the ordinary or extended stabilizer, and hence the precise identification of the target with a localized sphere, is not stated in the excerpt. Continuous cohomology with trivial coefficients must also be distinguished from the displayed coefficient action; constants may map into $(E_n)_0$, but the induced cohomology map need not be assumed injective.

The source reports that the degree-one class $\zeta$ in continuous stabilizer cohomology survives to a homotopy class of degree $-1$, using Hopkins--Miller and Devinatz--Hopkins fixed-point theory. It says large-prime collapse made this easy in that range but small primes had required these developments; no numerical prime bound or complete spectral-sequence hypothesis is supplied.

Morava K- and E-theory item 9 asks to find analogous classes in continuous degrees
\[
 3,5,\ldots,2n-1
\]
with the intended trivial coefficients, and determine whether they survive to nonzero homotopy classes in degrees $-3,-5,\ldots,-2n+1$ at all primes. Under the customary trivial-$\Z_p$ reading these would be classes $\zeta_{2k-1}\in H_c^{2k-1}(S;\Z_p)$ for $2\leq k\leq n$, with nonzero permanent-cycle images in coefficient degree $t=0$. The source presents their expected existence/survival as part of the chromatic-splitting picture; it does not ask for, or establish, all products or a full splitting theorem.

The registry reports the $\zeta$ case settled and the higher odd-degree survival not generally known. Its additional references to the image of $J$, level structures, cardinality arguments and height-two computations are not precise constructions or proofs, and its final progress sentence is truncated.
\subsection{Short English statement}
Find the higher odd continuous-cohomology classes predicted to produce degrees $-3,-5,\ldots,-2n+1$ in chromatic homotopy, and determine whether they are nonzero permanent cycles at every prime.
\subsection{Sources}
\begin{itemize}
\item Local registry: \url{data/unsolvedmath/problems.json}, permanent UnsolvedMath ID $11100019$, exact source code \texttt{catalogue entry}, attributed to Mark Hovey.
\item \emph{Hovey's Algebraic Topology Problem List}, Morava K- and E-theory item 9: \url{https://www-users.cse.umn.edu/~tlawson/hovey/}. Supplied access date: 2026-07-29; extraction: html-ordered-list; 65 top-level items in nine category pages.
\item Source attributions: Hopkins--Miller theorem and Devinatz--Hopkins fixed-point theory. The registry additionally names Morava and Ravenel, swindle/transfer methods, height-two cohomology and work at $p=3$, but supplies no exact papers or hypotheses; its final sentence is truncated.
\end{itemize}
\end{document}
